\documentclass[11pt]{article}
\usepackage{fontspec}
\usepackage[french]{babel}
\usepackage{amsmath}
\usepackage{amssymb}
\usepackage{amsthm}
\usepackage{mathrsfs}
\usepackage{tikz-cd}
\usepackage{enumitem}
\usepackage{lmodern}

\newcommand{\printedpage}[1]{\marginpar{\footnotesize p.~#1}}
\newcommand{\msp}[1]{\marginpar{\footnotesize [#1]}}

% short macros for printed notation
\newcommand{\uGamma}{\underline{\Gamma}}
\newcommand{\uFl}{\underline{\mathrm{Fl}}}
\newcommand{\uIs}{\underline{\mathrm{Is}}}
\newcommand{\uHom}{\underline{\mathrm{Hom}}}
\newcommand{\Ob}{\mathrm{Ob}\,}
\newcommand{\Fl}{\mathrm{Fl}\,}
\newcommand{\id}{\mathrm{id}}
\newcommand{\pr}{\mathrm{pr}}
\newcommand{\diag}{\mathrm{diag}\,}
\newcommand{\cardred}{\operatorname{card\,r\acute{e}d}}
\newcommand{\cocart}{\mathrm{cocart}}
\newcommand{\cart}{\mathrm{cart}}
\newcommand{\cc}{\mathrm{c}}
\newcommand{\circled}[1]{\tikz[baseline=(c.base)]\node[draw,circle,inner sep=0.8pt](c){\upshape #1};}

\begin{document}

\printedpage{149}
% head: A. Grothendieck : Dérivateurs, Ch. XVIII, Édition des Universités de Montpellier II et Paris VII

\msp{page 393}
Toutes ces assertions se prouvent facilement dans l’ordre où elles sont données, en utilisant 4.14.17 et les résultats généraux déjà prouvés. Bien sûr, on peut commencer par remplacer $\mathcal{C}$ par une sous-catégorie réduite $\mathcal{C}_0$ équivalente et $\mathcal{E}$ par $\mathcal{E}_0 = \mathcal{E}|\mathcal{C}$ \texttt{[plutôt} $\mathcal{E}|\mathcal{C}_0$\texttt{]}, de façon à se réduire (dans b) et c)) au cas où $\operatorname{card}\mathcal{C} \leqslant \pi$.

\textbf{Corollaire 4.14.19.} \textit{Sous les conditions de 4.14.17, la sous-catégorie $\uGamma_{\cocart}(\mathcal{E}/\mathcal{C})$ de $\mathcal{M} = \uGamma(\mathcal{E}/\mathcal{C})$ (cas $p$ cofibrant), resp. la sous-catégorie $\uGamma_{\cart}(\mathcal{E}/\mathcal{C})$ (cas $p$ fibrant) est accessible dans $\mathcal{M}$, donc accessible et le foncteur d’inclusion} \texttt{[est]} \textit{accessible.}

(\textbf{N.B.} Ces deux catégories sont strictement pleines dans $\uGamma(\mathcal{E}/\mathcal{C})$, donc l’assertion qu’elles soient accessibles dans celle-ci \textit{équivaut} à la dernière assertion.)

\textsc{Démonstration de} 4.14.19. On peut supposer $\mathcal{C}$ réduite, donc petite. Considérons le cas où $p$ est cofibrant. Pour toute flèche $\lambda : s \longrightarrow t$ dans $\mathcal{C}$, considérons le foncteur
\[
\varphi_\lambda \ : \ \mathcal{M} \longrightarrow \uFl(\mathcal{E}_t)
\]
défini par
\msp{page 394}
\[
\varphi_\lambda(f) \ = \ (\lambda_*(f(s)) \xrightarrow{f(\lambda)^\natural} f(t)) \ .
\]
C’est un foncteur accessible, car son composé avec les foncteurs source et but sont les foncteurs $\lambda_*\varepsilon_s$ et $\varepsilon_t$, donc accessibles, $\varepsilon_s$, $\varepsilon_t$ et $\lambda_*$ étant accessibles. Considérons la sous-catégorie strictement pleine
\[
\uIs(\mathcal{E}_t) \ \subset \ \uFl(\mathcal{E}_t) \ ,
\]
on sait qu’elle est accessible dans $\uFl(\mathcal{E}_t)$ (4.14.12). Donc son image inverse par $\varphi_\lambda$ est accessible dans $\mathcal{M}$, et par suite aussi
\[
\uGamma_{\cart}(\mathcal{E}/\mathcal{C}) \ \overset{\text{déf}}{=} \ \bigcap_{\lambda\in\Ob\mathcal{C}} \varphi_\lambda^{-1}(\uIs(\mathcal{E}_t))
\]
% [sic] \uGamma_{\cart} -> \uGamma_{\cocart}, the case treated is p cofibrant
% [sic] \lambda\in\Ob\mathcal{C} -> \lambda\in\Fl\mathcal{C}, \lambda runs over arrows of \mathcal{C}
comme intersection d’une \textit{petite} famille de sous-catégories accessibles dans $\mathcal{M}$, q.e.d. Le cas où $p$ est fibrant se traite de façon similaire, en introduisant les foncteurs
\[
\psi_\lambda \ : \ \mathcal{M} = \uGamma(\mathcal{E}/\mathcal{C}) \longrightarrow \uFl(\mathcal{E}_s)
\]
donnés par
\[
\psi_\lambda(f) \ = \ (f(s) \xrightarrow{f(\lambda)^\flat} \lambda^* f(t)) \ .
\]

\msp{page 395}
\textbf{Corollaire 4.14.20.} \textit{Soit $\pi_0$ un cardinal infini choisi comme dans 4.14.18 a).}
\begin{itemize}[leftmargin=2.5em]
\item[\circled{a}] \textit{Alors $\mathcal{M}_{\cc} = \uGamma_{\cocart}(\mathcal{E}/\mathcal{C})$ (resp. $\mathcal{M}_{\cc} = \uGamma_{\cart}$) est $L_{\pi_0}$-stable dans $\mathcal{M} = \uGamma(\mathcal{E}/\mathcal{C})$, ou ce qui revient au même (comme elle est strictement pleine), elle satisfait $L_{\pi_0}$, et le foncteur d’inclusion satisfait à $L_{\pi_0}$. Les $\varinjlim$ grandes devant $\pi_0$ dans cette catégorie se calculent argument par argument.}
\item[\circled{b}] \textit{Pour tout cardinal $\pi \geqslant \pi_0$, tout objet de $\mathcal{M}_{\cc}$ qui est $\pi$-accessible dans $\mathcal{M}$ l’est (tautologiquement) dans $\mathcal{M}_{\cc}$. Si $\pi \geqslant \sup(\pi_0, \cardred\mathcal{C})$, alors (par 4.14.18 b) tout objet $f$ de $\mathcal{M}_{\cc}$ tel que $f(s) \in \Ob\mathcal{E}_{s\,\pi}$ $\forall\ s \in S$} \texttt{[plutôt} $s \in \Ob C$\texttt{]} \textit{est dans $\mathcal{M}_\pi$} \textup{et} \textit{dans $\mathcal{M}_{\cc\,\pi}$.}
% foot: Version 17 décembre 2025

\printedpage{150}
% head: A. Grothendieck : Dérivateurs, Ch. XVIII, Édition des Universités de Montpellier II et Paris VII

\item[\circled{c}] \textit{Supposons $\pi_0$ choisi comme dans 4.14.18 c, et que de plus $\mathcal{M}_{\cc}$ soit $\pi_0$-accessible dans $\mathcal{M}$. (Un tel $\pi_0$ existe.) On a donc, pour $\pi \geqslant \pi_1 = \sup(\pi_0, \cardred\mathcal{C}, \cardred\mathcal{M}_\pi)$ tel que $\pi^{\pi_0} = \pi$,}
\[
(\mathcal{M}_{\cc})_\pi \ = \ \mathcal{M}_\pi{\cap}\mathcal{M}_{\cc} \ = \quad
\begin{array}{l}
\textit{Sous-catégorie strictement pleine de $\mathcal{M}_{\cc}$ formée des $f$}\\
\textit{telles que $f(s) \in \Ob\mathcal{E}_{s\,\pi}$ $\forall\ s \in \Ob\mathcal{C}$.}
\end{array}
\]

\msp{page 396}
\textit{Si de plus $\pi$ est choisi tel qu’on ait $\pi \geqslant \sup_{s\in\Ob\mathcal{C}} \cardred\mathcal{E}_{s\,\pi_0}$, alors}
\[
\left\{
\begin{array}{rrcl}
 & \mathcal{M}_{\cc\,\pi} & \xleftarrow{\ \sim\ } & \uGamma_{\cocart}(\mathcal{E}(\pi)/\mathcal{C})\\[1ex]
\textit{resp.} & \mathcal{M}_{\cc\,\pi} & \xleftarrow{\ \sim\ } & \uGamma_{\cart}(\mathcal{E}(\pi)/\mathcal{C}) \ ,
\end{array}
\right.
\]
\textit{donc ($\pi$ étant utile pour $\mathcal{M}_{\cc}$) une équivalence de catégories}
\[
\begin{array}{rrcl}
 & \uGamma_{\cocart}(\mathcal{E}/\mathcal{C}) & \xleftarrow{\ \sim\ } & \mathrm{Ind}_\pi(\uGamma_{\cocart}(\mathcal{E}(\pi)/\mathcal{C}))\\
\textit{resp.}\quad & \uGamma_{\cart}(\mathcal{E}/\mathcal{C}) & \xleftarrow{\ \sim\ } & \mathrm{Ind}_\pi(\uGamma_{\cart}(\mathcal{E}(\pi)/\mathcal{C})).
\end{array}
\]
\end{itemize}

\textsc{Démonstration}. Tout est clair dans l’ordre où c’est donné en utilisant les deux corollaires précédents.

\textbf{Corollaire 4.14.21.} \textit{Considérons un diagramme de foncteurs}
\[
\begin{tikzcd}[column sep=small, row sep=small]
\mathcal{E}_1 \arrow[dr, "p_1"] & & \mathcal{E}_2 \arrow[dl, "p_2"'] \\
& \mathcal{E}_0 \ , &
\end{tikzcd}
\tag*{$(*)$}
\]
\textit{les catégories et les foncteurs étant accessibles.}

\msp{page 397}
\begin{itemize}[leftmargin=2.5em]
\item[\textit{a)}] \textit{Le produit 2-fibré}
\[
\mathcal{M}_0 \ = \ \mathcal{E}_1 \overset{2}{\times}_{\mathcal{E}_0} \mathcal{E}_2
\]
\textit{est accessible, et les foncteurs canoniques $\mathcal{M} \longrightarrow \mathcal{E}_1$, $\mathcal{M} \longrightarrow \mathcal{E}_2$ sont accessibles.} % [sic] \mathcal{M} \longrightarrow \mathcal{E}_1, \mathcal{M} \longrightarrow \mathcal{E}_2 -> \mathcal{M}_0 \longrightarrow \mathcal{E}_1, \mathcal{M}_0 \longrightarrow \mathcal{E}_2, assertion a) concerns \mathcal{M}_0
\item[\textit{b)}] (\footnote[169]{Ceci est contenu dans 4.14.22.1 ci-dessus (p. 406), appliqué à $\Pi = \mathcal{E}_1{\times}\mathcal{E}_2$, $\Pi' = \mathcal{E}_0$, et $F, G : \Pi \rightrightarrows \Pi'$ donnés par $F = p_1\circ\pr_1$ et $G = p_2\circ\pr_2$. D’autre part, a) et c) résultent facilement de b).}). \textit{La catégorie $\mathcal{M}$ des triples $(x_1,x_2,u)$ avec $x_1 \in \Ob\mathcal{E}_1$, $x_2 \in \Ob\mathcal{E}_2$, et}
\[
u \ : \ p_1(x_1) \ \longrightarrow \ p_2(x_2) \ ,
\]
\textit{est accessible, et les foncteurs canoniques $\mathcal{M} \longrightarrow \mathcal{E}_1$, $\mathcal{M} \longrightarrow \mathcal{E}_2$ sont accessibles.}
\item[\textit{c)}] \textit{La sous-catégorie strictement pleine $\mathcal{M}_0$ de $\mathcal{M}$ est accessible dans $\mathcal{M}$.}
\end{itemize}

\noindent\textsc{Démonstration}. Soit $\mathcal{C}$ la catégorie
\[
\begin{tikzcd}[column sep=small]
1 \arrow[dr] & & 2 \arrow[dl] \\
& 0 \ , &
\end{tikzcd}
\]
et $\mathcal{E}$ la catégorie cofibrée sur $\mathcal{C}$ définie par le diagramme $(*)$. On est sous les conditions de 4.14.17. La catégorie $\Gamma = \uGamma(\mathcal{E}/\mathcal{C})$ s’interprète
% foot: Version 17 décembre 2025
\printedpage{151}
% head: A. Grothendieck : Dérivateurs, Ch. XVIII, Édition des Universités de Montpellier II et Paris VII
\msp{page 398}
comme la catégorie des quintuplets $(x_0,x_1,x_2,u_1,u_2)$, avec $x_i \in \Ob\mathcal{E}$ pour $i \in \{0,1,2\}$, et $u_1, u_2 \in \Fl\mathcal{E}_0$,
\[
\begin{tikzcd}[column sep=small, row sep=small]
p_1(x_1) \arrow[dr, "u_1"] & & p_2(x_2) \arrow[dl, "u_2"'] \\
& x_0 \ . &
\end{tikzcd}
\tag*{$(**)$}
\]
En vertu de 4.14.18, cette catégorie $\uGamma$ est accessible. La catégorie $\mathcal{M}$ est équivalente à la sous-catégorie strictement pleine $\Gamma_0$ de $\Gamma$, formée des quintuplets pour lesquels on a $u_2 \in \uIs\mathcal{E}_0$. On a en effet des foncteurs quasi-inverses l’un de l’autre % [sic] \Gamma_0 -> \Gamma_1, the subcategory is called \Gamma_1 in what follows
\[
\begin{tikzcd}
\mathcal{M} \arrow[r, shift left, "\alpha"] & \Gamma_1 \arrow[l, shift left, "\beta"]
\end{tikzcd}
\]
donnés par
\[
\begin{aligned}
\alpha(x_1,x_2,u) \ &= \ (p_2(x_2),x_1,x_2,u,\id_{p_2(x_2)})\\
\beta(x_0,x_1,x_2,u_1,u_2) \ &= \ (x_1,x_2,u_2^{-1}u_1) \ ,
\end{aligned}
\]
on a alors
\[
\begin{aligned}
\beta\alpha \quad &= \quad \id_{\mathcal{M}}\\
\alpha\beta \quad &\xrightarrow{\ \sim\ } \quad \id_{\Gamma_1} \ .
\end{aligned}
\]
D’autre part, $\Gamma_1$ est évidemment accessible dans $\mathcal{M}$, comme image inverse de $\uIs(\mathcal{E}_0)$ (qui % [sic] accessible dans \mathcal{M} -> accessible dans \Gamma, \Gamma_1 is a subcategory of \Gamma
\msp{page 399}
est accessible dans $\uFl(\mathcal{E}_0)$ par le foncteur accessible
\[
\uGamma \ \longrightarrow \ \uFl(\mathcal{E}_0)
\]
donné par $(x_1, x_1, x_2, u_1, u_2) \longmapsto u_2$. Cela prouve que $\mathcal{M}$, qui est équivalente à $\Gamma_1$, est également accessible, et l’accessibilité du foncteur canonique $\mathcal{M} \longrightarrow \mathcal{E}_i$ ($i \in \{1,2\}$) résulte de celle de $\Gamma \longrightarrow \mathcal{E}_i$, impliquant celle du composé $\Gamma_1 \hookrightarrow \Gamma \longrightarrow \mathcal{E}_i$. Cela prouve donc b). Pour prouver a), il suffit de voir que $\mathcal{M}_0$ est accessible dans $\mathcal{M}$, or c’est encore l’image inverse de $\uIs(\mathcal{E}_0) \subset \uFl(\mathcal{E}_0)$ par le foncteur accessible canonique % [sic] (x_1, x_1, x_2, u_1, u_2) -> (x_0, x_1, x_2, u_1, u_2), the quintuplets begin with x_0
\[
\mathcal{M} \ \longrightarrow \ \uFl(\mathcal{E}_0) \ ,
\]
donc $\mathcal{M}$ est bien accessible dans $\mathcal{M}$, comme image inverse d’une sous-catégorie accessible dans $\uFl(\mathcal{E})$. On peut d’ailleurs aussi interpréter $\mathcal{M}$, à équivalence de catégories près, comme % [sic] donc \mathcal{M} est bien accessible -> donc \mathcal{M}_0 est bien accessible, \mathcal{M}_0 is the subcategory in question
% [sic] \uFl(\mathcal{E}) -> \uFl(\mathcal{E}_0), as in the preceding display
% [sic] interpréter \mathcal{M} -> interpréter \mathcal{M}_0, as in the following display
\msp{page 400}
\[
\mathcal{M}_0 \ \simeq \ \uGamma_{\cocart}(\mathcal{E}/\mathcal{C}) \ ,
\]
les foncteurs canoniques $\mathcal{M}_0 \longrightarrow \mathcal{E}_i$ ($i \in \{0,1\}$) s’interprétant comme des foncteurs évaluation, ce qui donne a) directement, comme cas particulier de 4.14.19. % [sic] i \in \{0,1\} -> i \in \{1,2\}, the canonical functors of a) go to \mathcal{E}_1, \mathcal{E}_2
% foot: Version 17 décembre 2025

\printedpage{152}
% head: A. Grothendieck : Dérivateurs, Ch. XVIII, Édition des Universités de Montpellier II et Paris VII

\textbf{4.14.22.} \textsc{Démonstration de} 4.14.17. Le cas où $p$ est cofibrant a déjà été démontré (th. 4.10.8, p. 263). Nous allons voir que le cas où $p$ est fibrant se ramène au théorème précédent (dans un cas très particulier), en utilisant à fond les résultats généraux sur les parties accessibles d’une catégorie accessible. Notons que les assertions a) et b) du corollaire 4.14.18 se prouvent de façon AQT \texttt{[âne qui trotte]} sans avoir à faire appel à 4.14.17. Le a) implique déjà que les foncteurs évaluation $\varepsilon_s : \mathcal{M} \longrightarrow \mathcal{E}_s$ sont accessibles. Il reste donc à prouver
\msp{page 401}
que $\mathcal{M} = \uGamma(\mathcal{E}/\mathcal{C})$ est accessible, dans le cas où $p$ est \textit{fibrant}. On supposera de plus que $\mathcal{C}$ est petite (quitte à remplacer $\mathcal{C}$ par une petite sous-catégorie pleine équivalente).

Soient
\[
\begin{aligned}
\Pi \ &= \ \prod\nolimits_{s\in\Ob\mathcal{C}} \mathcal{E}_s\\
\Pi' \ &= \ \prod\nolimits_{\lambda\in\Fl\mathcal{C}} \mathcal{E}_{s(\lambda)} \ ,
\end{aligned}
\]
où $s(\lambda)$ désigne la \textit{source} $s$ de la flèche $\lambda : s \longrightarrow t$ dans $\mathcal{C}$. $\Pi$, $\Pi'$ sont des catégories accessibles, comme (petits) produits de catégories accessibles. (Cas particulier trivial de 4.10.8, pour une catégorie d’indices discrète.) On va définir deux foncteurs
\[
\Pi \ \overset{F}{\underset{G}{\rightrightarrows}} \ \Pi'
\]
en définissant pour tout $\lambda : s_0 \longrightarrow t_0$ dans $\Fl\mathcal{C}$ deux foncteurs
\msp{page 402}
\[
\left\{
\begin{array}{lcll}
\Pi & \overset{F_\lambda}{\underset{G_\lambda}{\rightrightarrows}} & \mathcal{E}_{s(\lambda)=s_0} & \\
F_\lambda & = & \pr_{s_0} & \text{projection de } \Pi = \prod\nolimits_s \mathcal{E}_s \text{ sur le facteur } \mathcal{E}_{s_0}\\
G_\lambda & = & \lambda^*\pr_{s_1} & \text{composé } \Pi \xrightarrow{\pr_{s_1}} \mathcal{E}_{s_1} \xrightarrow{\lambda^*} \mathcal{E}_{s_0}\\
 & & & \text{\texttt{[plutôt} } \lambda^*\pr_{t_0} \text{ \texttt{resp.} } \Pi \xrightarrow{\pr_{t_0}} \mathcal{E}_{t_0} \xrightarrow{\lambda^*} \mathcal{E}_{s_0}\text{\texttt{]}}
\end{array}
\right.
\qquad \bigl(\mbox{\footnotemark[170]}, \mbox{\footnotemark[171]}\bigr) \ .
\]
\footnotetext[170]{Il est immédiat que les $F_\lambda$, $G_\lambda$, donc $F$ et $G$, sont des foncteurs \textit{accessibles}.}
\footnotetext[171]{\textbf{N.B.} On suppose choisi pour tout $\lambda \in \Fl\mathcal{C}$ un foncteur image inverse $\lambda^* : \mathcal{E}_t \longrightarrow \mathcal{E}_s$. On suppose que si $\lambda$ est une flèche identique, $\id_s$, \texttt{[on]} choisit $\lambda^* = \id_{\mathcal{E}_s}$.}

Soit $f \in \mathcal{M} = \uGamma(\mathcal{E}/\mathcal{C})$, alors pour tout $\lambda \in \Fl\mathcal{C}$, $f$ définit une flèche
\[
\alpha_\lambda(f) \ : \quad \underbrace{f(s)}_{= F_\lambda(\varepsilon(f))} \quad \xrightarrow{f(\lambda)^\flat} \quad \underbrace{\lambda^*(f(t))}_{= G_\lambda(\varepsilon(f))} \ ,
\]
\textit{fonctorielle} en $f$. On a un foncteur
\[
\varepsilon \ = \ (\varepsilon_s)_{s\in\Ob\mathcal{C}} \ : \ \mathcal{M} \ \longrightarrow \ \Pi \ ,
\]
et pour tout $\lambda \in \Fl\mathcal{C}$ un homomorphisme fonctoriel
\[
\alpha_\lambda \ : \ F_\lambda \circ \varepsilon \ \longrightarrow \ G_\lambda \circ \varepsilon \ ,
\]
d’où un homomorphisme fonctoriel
\[
\alpha \ = \ (\alpha_\lambda)_{\lambda\in\Fl\mathcal{C}} \ : \ F \circ \varepsilon \ \longrightarrow \ G \circ \varepsilon \ .
\]
Soit alors $\mathcal{M}'$ la catégorie suivante :
% foot: Version 17 décembre 2025
\printedpage{153}
% head: A. Grothendieck : Dérivateurs, Ch. XVIII, Édition des Universités de Montpellier II et Paris VII
\msp{page 403}
\[
\mathcal{M}' \quad = \quad \parbox[t]{0.5\textwidth}{catégorie des couples $(X,u)$, avec $X$ dans $\Ob\Pi$, et $u : F(X) \longrightarrow G(X)$ une flèche dans $\Pi'$.}
\]
L’homomorphisme fonctoriel $\alpha : F \circ \varepsilon \longrightarrow G \circ \varepsilon$ définit un foncteur
\[
\mathcal{M} \ \longrightarrow \ \mathcal{M}' \ ,
\]
et il est immédiat que ce foncteur est pleinement fidèle. En fait, la catégorie $\mathcal{M}'$ est isomorphe à la catégorie formée des couples $(f_0,f_1)$ de deux applications
\[
\begin{array}{lcccccc}
f_0 : & \Ob\mathcal{C} & \longrightarrow & \Ob\mathcal{E} & s & \longmapsto & f_0(s)\\[1ex]
f_1 : & \Fl\mathcal{C} & \longrightarrow & \Fl\mathcal{E} & \lambda & \longmapsto & f_1(\lambda)
\end{array}
\]
de telle façon que $f_1$ et $f_0$ soient compatibles avec les applications source et but, i.e. qu’ils définissent un homomorphisme de \textit{diagrammes} $f : \diag\mathcal{C} \longrightarrow \diag\mathcal{E}$, et que cet homomorphisme soit inverse à droite de $\mathrm{diag}(p) : \diag\mathcal{E} \longrightarrow \diag\mathcal{C}$.

\msp{page 404}
Le foncteur précédent $\mathcal{M} \hookrightarrow \mathcal{M}'$ s’interprète comme le foncteur qui associe à toute sections $f$ de $\mathcal{E}$ sur $\mathcal{C}$, le morphisme de diagrammes correspondant. C’est donc l’inclusion d’une sous-catégorie \textit{strictement pleine} de $\mathcal{M}'$. Donc pour prouver que $\mathcal{M}$ est accessible, il suffit de prouver : % [sic] toute sections -> toute section, agreement of toute
\begin{itemize}[leftmargin=2.5em]
\item[a)] $\mathcal{M}'$ est accessible.
\item[b)] $\mathcal{M}$ est accessible dans $\mathcal{M}'$.
\end{itemize}
Prouvons d’abord b) (en admettant a). La sous-catégorie $\mathcal{M}$ de $\mathcal{M}'$ est définie par deux types des conditions, sur un objet $f = (f_0,f_1)$ de $\mathcal{M}'$ :
\begin{itemize}[leftmargin=2.5em]
\item[1°)] $\forall\ s \in \Ob\mathcal{C}$, $f_1(\id_s) = \id_{f_0(s)}$.
\item[2°)] Si $s \xrightarrow{\lambda} t \xrightarrow{\mu} u$ sont deux flèches composables dans $\mathcal{C}$, alors
\[
f_1(\mu\lambda) \ = \ f_1(\mu)f_1(\lambda) \ .
\]
\end{itemize}
Ces conditions s’explicitent l’une et \texttt{[l’autre par]} l’égalité d’une double-flèche dans $\mathcal{E}_s$,
\msp{page 405}
savoir
\[
f_0(s) \overset{f_1(\id_s)}{\underset{\id_{f_0(s)}}{\rightrightarrows}} f_0(s)
\]
dans le cas 1°), et
\[
f_0(s) \overset{f_1(\mu\lambda)^\flat}{\underset{(f_1(\mu)f_1(\lambda))^\flat}{\rightrightarrows}} (\mu\lambda)^*(f_0(u))
\]
dans le cas 2°). Dans l’un et l’autre cas, on a un foncteur
\[
\mathcal{M}' \ \xrightarrow{\gamma_s \text{ resp. } \gamma_{\lambda,\mu}} \ \uHom(I,\mathcal{E}_s)
\]
% foot: Version 17 décembre 2025

\end{document}
